\documentclass[10pt,a4paper]{amsart}
\usepackage[margin=19mm]{geometry}
\usepackage{amsmath,amssymb,mathtools}
\usepackage[scr=rsfs,cal=euler]{mathalfa}
\usepackage{xfrac}
\usepackage[unicode,hidelinks,bookmarks=false]{hyperref}
\newcommand{\ie}{i.e.\ }
\newcommand{\cf}{cf.\ }
\newcommand{\resp}{respectively\ }
\newcommand{\Subsection}{\subsection}
\providecommand{\nobreakdash}{\nobreak\hbox{-}}
\setlength{\emergencystretch}{2em}
\multlinegap0pt
\usepackage{bm}
\usepackage{mathtools}
\usepackage{amssymb}
\usepackage[shortlabels]{enumitem}
\usepackage{tikz}
\usepackage{xcolor}
\newtheorem{theorem}{Theorem}
\newtheorem{corollary}{Corollary}
\newtheorem{lemma}{Lemma}
\newcommand{\Aset}{\mathcal{A}}
\newcommand{\T}{\top}
\def\bA{\mathbf{A}}
\def\bB{\mathbf{B}}
\def\bI{\mathbf{I}}
\def\bPi{\boldsymbol{\Pi}}
\def\bR{\mathbf{R}}
\def\bS{\mathbf{S}}
\def\bW{\mathbf{W}}
\def\bX{\mathbf{X}}
\def\bbeta{\boldsymbol{\beta}}
\newcommand{\betahat}{\widehat{\boldsymbol\beta}}
\def\bu{\boldsymbol{u}}
\def\bv{\boldsymbol{v}}
\def\by{\boldsymbol{y}}
\DeclareMathOperator{\diag}{diag}
\newcommand{\trace}{\mathrm{trace}}
\title{\LARGE \textbf{Degrees of Freedom in Penalized Regression:\\ Model Selection with Adaptive Penalties}}
\author{Mauro Bernardi, Antonio Canale, Marco Stefanucci}
\date{Source: arXiv:2511.21595v3 (CC BY 4.0)}
\begin{document}
\maketitle

\noindent\emph{Source and licence.} \LARGE \textbf{Degrees of Freedom in Penalized Regression:\\ Model Selection with Adaptive Penalties}. Version arXiv:2511.21595v3 (CC BY 4.0), distributed by arXiv under the Creative Commons Attribution 4.0 International licence, http://creativecommons.org/licenses/by/4.0/, as recorded in the arXiv licence metadata for this exact version and shown on the abstract page https://arxiv.org/abs/2511.21595v3. Full licence notice: \texttt{licenses/arxiv-2511.21595-CC-BY-4.0.txt}.\par\smallskip

\noindent\textit{Editorial scope.} Counted unit: the article's theorem th:matrix\_Pi\_indefinite (source0.tex lines 822--825). The article's complete original proof of it is delivered with it (source0.tex lines 1325--1436, one \texttt{proof} environment of 112 lines, which the article places away from the statement). No mathematics is written, repaired or re-derived: the statement and the proof are the article's own lines, in the article's own notation, and no retained line is substituted or re-flowed.  Retained uncounted as the source-local context the proof needs: lines 753--756; lines 766--779; lines 782--796; lines 817--820; lines 1177--1180; lines 1731--1738; lines 1916--1932; lines 2040--2051.  Nothing was omitted from any retained span, so the package carries no omission record.  No retained line contains a citation, so the wrapper carries no bibliography and no bibliography entry.  No retained line contains a figure environment, an image-inclusion command or drawing code, so no image is delivered and the package declares no asset dependency.  Editorial formatting only, in addition: an AMS \texttt{amsart} preamble that restates the article's own declarations and macros used by the retained lines, namely the environments \texttt{theorem}, \texttt{corollary}, \texttt{lemma} on separate counters, together with the standard packages those lines need.  The retained environments print in this document as Theorem 1, Corollary 1, Lemma 1 in order of appearance, whereas the article prints the same items with the numbering of its own full document.  The complete native source of the article as distributed in the arXiv e-print for this exact version is bundled unchanged as \texttt{sources/189884/source0.tex}.
\begin{equation}
\label{eq:convex_regularized_problem_glasso}
\widehat{\boldsymbol{\beta}}=\arg\min_{\bbeta} \frac{1}{2}\Vert\by-\bX \bbeta \Vert_2^2+\gamma \sum_{g=1}^Gw_g\Vert \bbeta_g\Vert_2.
\end{equation}

\begin{theorem}
\label{th:df_glasso_ortho}
Let $\bX^{\T}\bX=\bI_p$, $\betahat$ be the solution to the Group Lasso problem in Equation \eqref{eq:convex_regularized_problem_glasso} with groups cardinality equal to $n_1, \ldots, n_G,$ %, weights $w_g=1 \ \forall g$
and $\gamma \in (\gamma_l, \gamma_{l+1})$. Denote with $\Aset_G$ the corresponding active group set. An unbiased estimate of the degrees of freedom is 
\begin{equation}
\label{eq:dof_glasso_ortho}
\widehat{df}_\gamma = \trace \big[ \big(\bI_n+\gamma \bB \big)^{-1} \bA  \big],
\end{equation} 
where $\bA = \bX_{\Aset_G}\bX_{\Aset_G}^{\T},$ $\bB = \bX_{\Aset_G}\bPi_{\Aset_G}\bX_{\Aset_G}^{\T}$, and 
\begin{equation}
\label{eq:ABPi_g_definition} 
\bPi_{\Aset_G} = \mathrm{blockdiag}\bigg(w_g\bigg[\frac{\bI_{n_g}}{\Vert\betahat_g\Vert_2}-\frac{\betahat_g\betahat_g^{\T}}{\Vert\betahat_g\Vert_2^3}\bigg]\bigg).
\end{equation}
\end{theorem}

\begin{theorem}
\label{th:df_glasso_nonortho} 
Let  $\bX^{\T}\bX \neq\bI_p$, $\betahat$ the solution to the Group Lasso problem in Equation \eqref{eq:convex_regularized_problem_glasso} with groups cardinality equal to $n_1, \ldots, n_G$ %, weights $w_g=1 \ \forall g$ 
and $\gamma \in (\gamma_l, \gamma_{l+1})$. Denote with $\Aset_G$ the corresponding active group set. An unbiased estimate of the degrees of freedom is  
\begin{equation}
\label{eq:dof_glasso_nonortho}
\widehat{df}_\gamma = \trace \big[ \big(\bI_n+\gamma \bB\big)^{-1} \bA  \big],
\end{equation}
where $\bA = \bX_{\Aset_G}(\bX_{\Aset_G}^{\T}\bX_{\Aset_G})^{-1}\bX_{\Aset_G}^{\T}$,
$\bB = \bX_{\Aset_G}(\bX_{\Aset_G}^{\T}\bX_{\Aset_G})^{-1} \bPi_{\Aset_G} (\bX_{\Aset_G}^{\T}\bX_{\Aset_G})^{-1} \bX_{\Aset_G}^{\T}$,
and
\begin{equation*}
\bPi_{\Aset_G} = \mathrm{blockdiag}\bigg(w_g\bigg[\frac{\bI_{n_g}}{\Vert\betahat_g\Vert_2}-\frac{\betahat_g\betahat_g^{\T}}{\Vert\betahat_g\Vert_2^3}\bigg]\bigg).
\end{equation*}
\end{theorem}

\begin{theorem}
\label{th:monotonicity_df_glasso}
If the matrix $\bPi_{\Aset_G} + \gamma \mathrm{d} \bPi_{\Aset_G} / \mathrm{d}\gamma $ is positive definite, then the estimated degrees of freedom of Group Lasso given by Theorem \ref{th:df_glasso_ortho} and \ref{th:df_glasso_nonortho} are non-increasing function of $\gamma$, for $\gamma \in (\gamma_l, \gamma_{l+1})$.
\end{theorem}

\begin{equation}
\label{eq:Pi_Aset_definition}
\bPi_g = w_g \Bigg(\frac{\bI_{p}}{\Vert\betahat_g\Vert_2}-\frac{\betahat_g\betahat_g^{\T}}{\Vert\betahat_g\Vert_2^3}\Bigg) \quad \text{and} \quad \bPi_{\Aset_G} = \mathrm{blockdiag}(\bPi_1, \ldots, \bPi_G),
\end{equation}

\begin{corollary}
\label{coro:dbeta_dgamma_orto}
Under the setting of Theorem \ref{th:monotonicity_df_glasso}, for the orthonormal design, $\bX^{\T}\bX=\bI_p$, we have:
\begin{equation*}
\frac{\mathrm{d}\betahat_g}{\mathrm{d}\gamma} = - \frac{w_g}{r_g}  \betahat_g, \qquad  \frac{{\mathrm d}r_{g}}{{\mathrm d} \gamma} = -w_g, \qquad \frac{{\mathrm d}\bPi_{g}}{{\mathrm d} \gamma} = \frac{w_g}{r_g} \bPi_{g} , \qquad \text{for each } g \in \Aset_G,
\end{equation*}
where $r_g= \Vert\betahat_g\Vert_2$, which means that the solution scales linearly in the direction of $\betahat_g$ as $\gamma$ increases, for each active group.
\end{corollary}

\begin{lemma}
\label{lemma:Pi_eigenvalues_all}
Let $\boldsymbol{u},\boldsymbol{v} \in \mathbb{R}^n$ with $n\geq 2$, $\bu\neq\bv$, $a,b,c\in\mathbb{R}\setminus \{0\}$, then the spectrum of the matrix:
\begin{equation*}
\boldsymbol{\Pi} = a \boldsymbol{I}_n + b \boldsymbol{u} \boldsymbol{u}^{\T} + c(\boldsymbol{u}\boldsymbol{v}^{\T} + \boldsymbol{v}\boldsymbol{u}^{\T})\in\mathbb{S}^n,
\end{equation*}
is
%
\begin{align*}
\lambda_1 &= \frac{1}{2}(2 a+b \Vert\bu\Vert_2^2+2 c \bv^{\T}\bu) + \frac{1}{2}\sqrt{\Delta}, \\
\lambda_2 &= a, \qquad \text{with multiplicity } n - 2, \\
\lambda_3 &= \frac{1}{2}(2 a+b \Vert\bu\Vert_2^2+2 c \bv^{\T}\bu) - \frac{1}{2}\sqrt{\Delta},
\end{align*}
and $\Delta =4c^2\Vert\bu\Vert_2^2\Vert\bv\Vert_2^2+b^2\Vert\bu\Vert_2^4+4bc\Vert\bu\Vert_2^2\bv^{\T}\bu$, if $n>2$ and $(\lambda_1,\lambda_3)$ if $n=2$.
%
%Guarda qua: $\lambda_1+\lambda_3=\trace(\bPi_{\mathcal{S}})=2 a+b \Vert\bu\Vert_2^2+2 c \bv^{\T}\bu$
\end{lemma}

\begin{lemma}
\label{lemma:Delta_sign}
Let $\boldsymbol{u},\boldsymbol{v} \in \mathbb{R}^n$ with $n\geq 2$, $\gamma, w, r \in \mathbb{R}_+$, then the quantity 
$$
\Delta=\frac{w^2}{\Vert \boldsymbol{u} \Vert_2^6} \left(
-3\gamma^2(\boldsymbol{u}^{\T} \boldsymbol{v})^2 
+ 2\gamma \Vert \boldsymbol{u} \Vert_2^2 (\boldsymbol{u}^{\T} \boldsymbol{v}) 
+ 4\gamma^2 \Vert \boldsymbol{u} \Vert_2^2 \Vert \boldsymbol{v} \Vert_2^2 
+ \Vert \boldsymbol{u} \Vert_2^4 \right)=\frac{w^2}{\Vert \boldsymbol{u} \Vert_2^6}\mathcal{Q}(\boldsymbol{u}^{\T}\boldsymbol{v}),
$$ 
is positive for all real vectors $\boldsymbol{u}, \boldsymbol{v}$, and it is strictly positive unless $\rho_{u,v} = -1$ and $\gamma \Vert \boldsymbol{v} \Vert_2 = \Vert \boldsymbol{u} \Vert_2$. 
\end{lemma}

\begin{theorem}
\label{th:matrix_Pi_indefinite}
    If $\bX^{\T} \bX = \bI_p$ or $\betahat_{\Aset_G}$ is an eigenvector of the matrix $(\bX_{\Aset_G}^{\T} \bX_{\Aset_G} + \gamma \bPi_{\Aset_G})^{-1} \bW_{\Aset_G}$, the matrix $\bPi_{\Aset_G} + \gamma \mathrm{d} \bPi_{\Aset_G} / \mathrm{d}\gamma$ is positive semidefinite. Otherwise, it is indefinite.
\end{theorem}

\begin{proof}[Proof of Theorem \ref{th:matrix_Pi_indefinite}]
Consider the derivative of $\bPi_{\Aset_G}$ wrt $\gamma$. Since $\bPi_{\Aset_G}$ is block-diagonal with blocks defined in Equation \eqref{eq:Pi_Aset_definition}, then the derivative can be computed blockwise:
\begin{equation*}
\frac{\mathrm{d}\bPi_{\Aset_G} }{\mathrm{d}\gamma}= \mathrm{blockdiag}\Big(\frac{\mathrm{d}\bPi_{1}}{\mathrm{d}\gamma}, \ldots, \frac{\mathrm{d}\bPi_{G}}{\mathrm{d}\gamma}\Big).
\end{equation*}
Applying the chain rule for the derivative, and defining $r_g=\|\betahat_{g}\|_2$, we have
\begin{equation}
\label{eq:Pi_derivative}
\frac{\mathrm{d}\bPi_{g}}{\mathrm{d}\gamma}
= -\frac{w_g}{r_g^2}\frac{\mathrm{d}r_g}{\mathrm{d}\gamma}\Big(\bI_{n_g}-\frac{3}{r_g^2} \widehat{\bbeta}_g \widehat{\bbeta}_g^{\T} \Big)-  \frac{w_g}{r_g^3} \Big( \frac{\mathrm{d}\betahat_g}{\mathrm{d}\gamma} \widehat{\bbeta}_g^{\T} + \widehat{\bbeta}_g \frac{\mathrm{d}\betahat_g}{\mathrm{d}\gamma}^{\T} \Big),\quad g \in \Aset_G,
\end{equation}
$\frac{\mathrm{d}\betahat_g}{\mathrm{d}\gamma}=\frac{\mathrm{d}\betahat_g}{\mathrm{d}\betahat_{\Aset_G}}\frac{\mathrm{d}\betahat_{\Aset_G}}{\mathrm{d}\gamma}=\bS_{g}\frac{\mathrm{d}\betahat_{\Aset_G}}{\mathrm{d}\gamma}$, and $\bS_{g}\in\{0,1\}^{n_g\times\vert\Aset_G\vert}$ denotes the selection matrix obtained from the identity matrix $\bI_{\vert\Aset_G\vert}$ by retaining only the rows corresponding to the $g$-th group. To compute the derivative $\frac{\mathrm{d}\widehat{\bbeta}_{\Aset_G}}{\mathrm{d}\gamma}$, we apply the implicit function theorem to the equation:
\begin{equation*}
F(\widehat{\bbeta}_{\Aset_G}, \gamma) =-\bX_{\Aset_G}^{\T}\left(\by-\bX_{\Aset_G}\betahat_{\Aset_G}\right)+\gamma\widetilde{\bbeta}_{\Aset_G} =\boldsymbol{0}.
\end{equation*}
Differentiating $F(\widehat{\bbeta}_{\Aset_G}, \gamma)$ with respect to $\gamma$ and solving for $\frac{\mathrm{d}\widehat{\bbeta}_{\Aset_G}}{\mathrm{d} \gamma}$ gives:
\begin{equation}
\label{eq:IFT_Pi}
\frac{\mathrm{d}\widehat{\bbeta}_{\Aset_G}}{\mathrm{d}\gamma} = -\left( \frac{\partial F}{\partial \widehat{\bbeta}_{\Aset_G}} \right)^{-1}  \frac{\partial F}{\partial \gamma}\in\mathbb{R}^{\vert\Aset_G\vert},
\end{equation}
where
\begin{equation*}
\frac{\partial F}{\partial \widehat{\bbeta}_{\Aset_G}} = \bX_{\Aset_G}^{\T} \bX_{\Aset_G}+ \gamma \bPi_{\Aset_G}\in\mathbb{S}^{\vert\Aset_G\vert}_+,
\quad
\frac{\partial F}{\partial \gamma} =\widetilde{\bbeta}_{\Aset_G}\in\mathbb{R}^{\vert\Aset_G\vert},
\end{equation*}
and, substituting in Equation \eqref{eq:IFT_Pi}, we get the final expression for the derivative:
\begin{equation}
\label{eq:IFT_Pi_FINAL}
\frac{\mathrm{d}\widehat{\bbeta}_{\Aset_G}}{\mathrm{d}\gamma}
= -\left( \bX_{\Aset_G}^{\T} \bX_{\Aset_G} + \gamma \bPi_{\Aset_G} \right)^{-1}\widetilde{\bbeta}_{\Aset_G}=-\bR_{\Aset_G}\widetilde{\bbeta}_{\Aset_G}=-\bR_{\Aset_G}^\star\widehat{\bbeta}_{\Aset_G}\in\mathbb{R}^{\vert\Aset_G\vert}.
\end{equation}
where $\bR_{\Aset_G}^\star=\bR_{\Aset_G}\bW_{\Aset_G},\bW_{\Aset_G}=\diag\Big\{\frac{w_g}{\Vert\widehat{\bbeta}_{g}\Vert_2}\boldsymbol{\iota}_{n_g}\Big\}_{g\in\Aset_G}$. Moreover, letting $r_{\Aset_G} = \Vert \widehat{\bbeta}_{\Aset_G}\Vert_2$ and $r_g= \Vert \widehat{\bbeta}_{g}\Vert_2$, which are functions of $\gamma$, we aim to compute their derivative with respect to $\gamma$. Using the chain rule, we obtain the following:
\begin{equation*}
%\label{eq:drdgamma_step1}
\frac{\mathrm{d}r_{\Aset_G}}{\mathrm{d}\gamma}
= \frac{\mathrm{d}}{\mathrm{d}\gamma} \left( \Vert \widehat{\bbeta}_{\Aset_G} \Vert_2 \right)
= \frac{1}{\Vert \widehat{\bbeta}_{\Aset_G} \Vert_2}\widehat{\bbeta}_{\Aset_G}^{\T} \frac{\mathrm{d}\widehat{\bbeta}_{\Aset_G}}{\mathrm{d}\gamma}.
\end{equation*}
Substituting the expression for \( \frac{\mathrm{d}\widehat{\bbeta}_{\Aset_G}}{\mathrm{d}\gamma} \) from Equation~\eqref{eq:IFT_Pi_FINAL}, we get:
\begin{equation*}
\frac{\mathrm{d}r_{\Aset_G}}{\mathrm{d}\gamma}
= \frac{1}{\Vert \widehat{\bbeta}_{\Aset_G} \Vert_2}
\left( -\widehat{\bbeta}_{\Aset_G}^{\T} \bR_{\Aset_G} \widetilde{\bbeta}_{\Aset_G} \right)=-\frac{ \widehat{\bbeta}_{\Aset_G}^{\T} \bR_{\Aset_G} \widetilde{\bbeta}_{\Aset_G} }{ \Vert \widehat{\bbeta}_{\Aset_G} \Vert_2}=-\frac{ \widehat{\bbeta}_{\Aset_G}^{\T} \bR_{\Aset_G}^\star \widehat{\bbeta}_{\Aset_G} }{\Vert \widehat{\bbeta}_{\Aset_G} \Vert_2},
\end{equation*}
and
\begin{equation}
\label{eq:dr_gdgamma}
\frac{\mathrm{d}r_g}{\mathrm{d}\gamma}
= \frac{\mathrm{d}}{\mathrm{d}\gamma} \left( \Vert \widehat{\bbeta}_{g} \Vert_2 \right)
= \frac{1}{\Vert \widehat{\bbeta}_{g} \Vert_2}\widehat{\bbeta}_{g}^{\T} \frac{\mathrm{d}\widehat{\bbeta}_{g}}{\mathrm{d}\gamma}=-\frac{1}{\Vert \widehat{\bbeta}_{g} \Vert_2}\widehat{\bbeta}_{g}^{\T} \Big[\bR_{\Aset_G}^\star\betahat_{\Aset_G}\Big]_g.%=-\frac{1}{r_g}\bu^{\T}\bv.
\end{equation}
Substituting the expression for $\frac{\mathrm{d}\betahat_g}{\mathrm{d}\gamma}$ into Equation 
\eqref{eq:Pi_derivative}, we get:
\begin{equation}
\label{eq:Pi_derivative_final}
\frac{\mathrm{d}\bPi_{g}}{\mathrm{d}\gamma}
= -\frac{w_g}{r_g^2}\frac{\mathrm{d}r_g}{\mathrm{d}\gamma}\Big(\bI_{n_g}-\frac{3}{r_g^2} \widehat{\bbeta}_g \widehat{\bbeta}_g^{\T} \Big)+  \frac{w_g}{r_g^3} \Big(\Big[\bR_{\Aset_G}^\star\widehat{\bbeta}_{\Aset_G}\Big]_g \widehat{\bbeta}_g^{\T} + \widehat{\bbeta}_g \Big[\bR_{\Aset_G}^\star\widehat{\bbeta}_{\Aset_G}\Big]_g^{\T} \Big).
\end{equation}
We are interested in finding all the eigenvalues of the matrix $\bPi_{\Aset_G}+\gamma \mathrm{d}  \bPi_{\Aset_G}/\mathrm{d} \gamma$. Let $\bu_g = \widehat{\bbeta}_g,\quad \bv_g = \left[\bR_{\Aset_G}^\star \widehat{\bbeta}_{\Aset_G} \right]_g$, the eigenvalues of the $g$-th block of $\bPi_{\Aset_G}+\gamma \mathrm{d}  \bPi_{\Aset_G}/\mathrm{d} \gamma$ can be calculated using Lemma \ref{lemma:Pi_eigenvalues_all} in the Supplementary Materials. Specifically
\begin{align*}
\lambda_{1,g} &= \frac{1}{2}(2 a_g+b_g \Vert\bu_g\Vert_2^2+2 c_g \bv_g^{\T}\bu_g) + \frac{1}{2}\sqrt{\Delta_g}, \\
\lambda_{2,g} &= a_g, \qquad \text{with multiplicity } n_g - 2, \\
\lambda_{3,g} &= \frac{1}{2}(2 a_g+b_g \Vert\bu_g\Vert_2^2+2 c_g \bv_g^{\T}\bu_g) - \frac{1}{2}\sqrt{\Delta_g},
\end{align*}
where
\begin{equation}
\label{eq:abc_definition}
a_g=\gamma\frac{w_g}{r_g^3}\bu_g^{\T}\bv_g+\frac{w_g}{r_g},\qquad
b_g=-\Big(\gamma\frac{3w_g}{r_g^5}\bu_g^{\T}\bv_g+\frac{w_g}{r_g^3}\Big),\qquad
c_g=\gamma\frac{w_g}{r_g^3}>0,
\end{equation}
and
\begin{equation}
\label{eq:Pi_eigen_Delta}
\Delta_g=4c_g^2r_g^2\Vert\bv_g\Vert_2^2+b_g^2r_g^4+4b_gc_gr_g^2\bu_g^{\T}\bv_g=-\frac{3\gamma^2w_g^2}{r_g^6}(\bu_g^{\T}\bv_g)^2+\frac{2\gamma w_g^2}{r_g^4}\bu_g^{\T}\bv_g+\frac{4\gamma^2 w_g^2\Vert\bv_g\Vert_2^2}{r_g^4}+\frac{w_g^2}{r_g^2},
\end{equation}
for $g=1,\dots,G$. Moreover, $\Delta_g\geq0$, by Lemma \ref{lemma:Delta_sign} in the Supplementary Materials. Let us now consider the sum and product of the eigenvalues $\lambda_{1,g}$ and $\lambda_{3,g}$. We have:
\begin{equation}
\label{eq:lambda1_plus_lambda3_is_lambda2}
\begin{aligned}
\lambda_{1,g}+\lambda_{3,g}&=2 a_g+b_g \Vert\bu_g\Vert_2^2+2 c_g \bv_g^{\T}\bu_g\\
&=2\Big(\gamma\frac{w_g}{r_g^3}\bu_g^{\T}\bv_g+\frac{w_g}{r_g}\Big)-\Big(\gamma\frac{3w_g}{r_g^5}\bu_g^{\T}\bv_g+\frac{w_g}{r_g^3}\Big)r_g^2+2\gamma\frac{w_g}{r_g^3}\bu_g^{\T}\bv_g\\
&=\frac{\gamma w_g}{r_g^3}\bu_g^{\T}\bv_g+\frac{w_g}{r_g}=a_g=\lambda_{2,g},
\end{aligned}
\end{equation}
therefore $\lambda_{1,g}+\lambda_{3,g}\geq0$ if $\rho_{u,v}\geq-\frac{r_g}{\gamma\Vert\bv_g\Vert_2}$, where $\rho_{u,v}=\frac{\bu_g^{\T}\bv_g}{\Vert\bu_g\Vert_2\Vert\bv_g\Vert_2}$. As concerns the product of $\lambda_{1,g}$ and $\lambda_{3,g}$, we have:
\begin{equation*}
\begin{aligned}
\lambda_{1,g}\lambda_{3,g}&=\frac{1}{4}(2 a_g+b_g \Vert\bu_g\Vert_2^2+2 c_g \bv_g^{\T}\bu_g)^2 - \frac{1}{4}\Delta_g\\
&=a_g^2+\frac{1}{4}b_g^2r_g^4+c_g^2(\bu_g^{\T}\bv_g)^2+a_gb_gr_g^2+2a_gc_g(\bu_g^{\T}\bv_g)+b_gc_gr_g^2(\bu_g^{\T}\bv_g)\\
&\qquad-c_g^2r_g^2\Vert\bv_g\vert_2^2-\frac{1}{4}b_g^2r_g^4-b_gc_gr_g^2\bu_g^{\T}\bv_g\\
&= a_g^2+a_gb_gr_g^2-c_g^2\Vert\bu_g\Vert_2^2\Vert\bv_g\Vert_2^2(1-\rho_{u,v}^2)+2a_gc_g\bu_g^{\T}\bv_g.
\end{aligned}
\end{equation*}
Now, we substitute the values of $a_g$, $b_g$ and $c_g$ in Equation \eqref{eq:abc_definition} and we get:
\begin{equation*}
\begin{aligned}
\lambda_{1,g}\lambda_{3,g}&=\Big(\gamma\frac{w_g}{r_g^3}\bu_g^{\T}\bv_g+\frac{w_g}{r_g}\Big)^2-\Big(\gamma\frac{w_g}{r_g^3}\bu_g^{\T}\bv_g+\frac{w_g}{r_g}\Big)\Big(\gamma\frac{3w_g}{r_g^5}\bu_g^{\T}\bv_g+\frac{w_g}{r_g^3}\Big)r_g^2\\
&\qquad-\gamma^2\frac{w_g^2}{r_g^6}r_g^2\Vert\bv_g\Vert_2^2(1-\rho_{u,v}^2)+2\Big(\gamma\frac{w_g}{r_g^3}\bu_g^{\T}\bv_g+\frac{w_g}{r_g}\Big)\gamma\frac{w_g}{r_g^3}\bu_g^{\T}\bv_g\\
&=\frac{\gamma^2w_g^2}{r_g^6}\Vert\bu_g\Vert_2^2\Vert\bv_g\Vert_2^2\big(\rho_{u,v}^2-1\big).
\end{aligned}
\end{equation*}
Therefore, $\lambda_{1,g}\lambda_{3,g}<0$ and the matrix $\bPi_{\Aset_G}+\gamma\mathrm{d}\bPi_{\Aset_G}/\mathrm{d}\gamma$ is indefinite, unless $\rho_{u,v}=1$ (e.g. $\bu_g\propto\bv_g$) when it is equal to $\lambda_{1,g}\lambda_{3,g}=0$. In particular, when $\bX^{\T}\bX=\bI_p$ then $\rho_{u,v}=1$ (see Corollary \ref{coro:dbeta_dgamma_orto} in Supplementary Materials) and $\lambda_{3,g}=\frac{1}{2}(2a_g+b_g r_g^2+2cv\bu_g^{\T}\bv_g)-\frac{1}{2} \sqrt{\Delta_g}$ where $\Delta_g$ is defined in Equation \eqref{eq:Pi_eigen_Delta} as:
\begin{equation*}
\begin{aligned}
\Delta_{g,\textrm{ortho}}&=-\frac{3\gamma^2w_g^2}{r_g^6}\rho_{u,v}^2r_g^2\Vert\bv_g\Vert_2^2+\frac{2\gamma w_g^2}{r_g^4}\rho_{u,v} r_g\Vert\bv_g\Vert_2+\frac{4\gamma^2w_g^2\Vert\bv_g\Vert_2^2}{r_g^4}+\frac{w_g^2}{r_g^2}\\
&=\frac{\gamma^2w_g^2\Vert\bv_g\Vert_2^2}{r_g^4}+\frac{2\gamma w_g^2\Vert\bv_g\Vert_2}{r_g^3}+\frac{w_g^2}{r_g^2}=\Bigg(\frac{\gamma w_g\Vert\bv_g\Vert_2}{r_g^2}+\frac{w_g}{r_g}\Bigg)^2=a_g^2.
\end{aligned}
\end{equation*}
Moreover, by Equation \eqref{eq:lambda1_plus_lambda3_is_lambda2}, the first part of $\lambda_{3,g}$, $2a_g+b_g r_g^2+2c_g\bu_g^{\T}\bv_g=a_g$, therefore $\lambda_{3,g}=a-\sqrt{a^2}=0$ and $\lambda_{1,g}=\lambda_{2,g}$. Also, $\bu_g^{\T}\bv_g=\Vert\bu_g\Vert_2\Vert\bv_g\Vert_2$ and $a_g$ defined in Equation \eqref{eq:abc_definition} is strictly positive, and the matrix $\bPi_{g} + \gamma \mathrm{d}\bPi_{g}/\mathrm{d}\gamma$ is positive semidefinite. To check if the orthonormal design is the only setting leading to $\rho_{u,v}=1$, we observe that $\rho_{u,v}=1$ implies $\bR_{\Aset_G}^\star\widehat{\bbeta}_{\Aset_G} = c \widehat{\bbeta}_{\Aset_G}$, which in turn means that $\widehat{\bbeta}_{\Aset_G}$ is an eigenvector of the matrix $\bR^*_{\Aset_G} = \bR_{\Aset_G}\bW_{\Aset_G} = (\bX_{\Aset_G}^{\T} \bX_{\Aset_G} + \gamma \bPi_{\Aset_G})^{-1} \bW_{\Aset_G}$, which completes the proof.
\end{proof}

\end{document}
